% Page 058

\seite{58}

\margmark{16. Juni.}%
In der heutigen Schulvorstandssitzung wurde der\ob
als erster Ferientag für die Bauferien der 17.\ Juni\ob
festgelegt, indesfolge die allgemeinen Prüfungen die\ob
Ferienzeit in diesem Jahre schon früh ihren Anfang\ob
nehmen mußte.\ob
Zum Vorsitzenden des hiesigen Schulvorstandes wurde\ob
Lehrer Killipp vorgeschlagen.

In der heutigen Gemeinderatssitzung wurde\ob
dem Lehrer auf dessen Antrag das Gemeinde-Nutzungsrecht\ob
gewährt.

\margmark{4. Juli.}%
Die Heuferien wurden am 3.\ und 4.\ d. M.\ infolge\ob
mangelhaften Heuwuchses unterbrochen.

Das laufende Jahr scheint allgemein schlecht auszufallen.\ob
Der Getreidestand blieb den vorjährigen gegenüber\ob
weit zurück; der Haber steht schlecht. Überall jammert\ob
man nach Regen. Betrachten wir Wiesen, Äcker\ob
und Gärten — es sieht schlecht aus! Kein Regen,\ob
allgemeine Wasserarmut, überall große Dürre. Womit\ob
soll der Bauer sein Vieh füttern? Er sieht sich gezwungen,\ob
manches Stücklein im Herbst billig zu verkaufen, wo sich die\ob
Vorsehung nicht bald erbarmet und nicht hilft. „Der\ob
Mensch denkt und Gott lenkt!"

\margmark{31. Juli.}%
Mit unverminderter Stärke hält die fast unerträgliche\ob
Hitze an, der letzte Rest der Lebensfähigkeit ist längst\ob
dahingeschwunden. Waldbrände sind an der Tagesordnung,\ob
besonders in den angrenzenden Nachbarbezirken — in dem\ob
Kyllburgerwald, wo es bereits viermal im Laufe dieses\ob
Sommers brannte. Kein grüner Halm ist mehr auf den Wiesen\ob
zu finden. Gartenfrüchte gibt es nicht, wo man nicht mit\ob
unermüdlichem Fleiße Wasser schleppt, aber auch dort ist der\ob
verdiente Lohn kaum entsprechend. Oft und zu erscheinen\ob
Gewitter — doch ohne Regen. 38°C Wärme wurden\ob
mehrere Tage festgestellt. Wann kommt die längst\ob
ersehnte Linderung?!

\margmark{20. August.}%
Am 18.\ d. Mts. machte die hiesige Schule ihren Jahresausflug\ob
nach Thülecken und dem Wasserfall bei Scheid. Abends\ob
Heimkehr. Abends wurde gesammelt für wohlthätige\ob
Zwecke, wo von den Kindern an freiwilligen Gaben geopfert:\ob
für den Auslandsdeutschen Verein 9 M., zur Erhaltung der\ob
Kriegergräber 7 M., zum Schuljubiläum 5 M.,\ob
und für die Diaspora-Kollekte 2 M.\ 6 Pf.
\pend
